%%
%% The "title" command has an optional parameter,
%% allowing the author to define a "short title" to be used in page headers.
\title{LAB 4:  CST Element and FEM Assembly}
% \date{October 7, 2026}

%%
\author{Adhilsha A (SE25PCSE007)}
\maketitle

\tableofcontents

\section{Single CST Element}

The objective is to  calculate the element area, strain-displacement matrix,
material constitutive matrix, element stiffness matrix, nodal displacements,
strains, and stresses for a triangular Constant Strain Triangle element using the finite
element method. We will also calculate the support reactions and compare the
computed response with the expected analytical behaviour. The system is given as:
\begin{table}[h!]
\centering
\begin{tabular}{ll}
\toprule
\textbf{Quantity} & \textbf{Value} \\
\midrule
Element type & Constant Strain Triangle (CST) \\
Node 1 & $(0,0)\ \si{\milli\metre}$ \\
Node 2 & $(1,0)\ \si{\milli\metre}$ \\
Node 3 & $(0,1)\ \si{\milli\metre}$ \\
Young's modulus, $E$ & $\SI{200000}{\mega\pascal}$ \\
Poisson's ratio, $\nu$ & $0.30$ \\
Thickness, $t$ & $\SI{1}{\milli\metre}$ \\
Analysis assumption & Linear elastic plane stress \\
\bottomrule
\end{tabular}
\end{table}

\subsection{Code}

The matlab code along with the comments are added below:
\begin{lstlisting}[style=matlabcode]
%% --- MATERIAL PROPERTIES --- %%
E = 200000;  % MPa
t = 1;  % mm
nu = 0.30; 

% node coordinates and element numbering
nodes = [0 0; 1 0; 0 1];
elements = [1 2 3];

% construct CST matrices
[A1, B1, D1, Ke1] = cst_matrices(nodes(elements,:), E, nu, t);

% boundary conditions
F1 = [0; 0; 0; 0; 0; -100];
fixedDofs1 = [1 2 4];

% solving and results
U1 = solve_with_constraints(Ke1, F1, fixedDofs1);
R1 = Ke1*U1 - F1;
strain1 = B1*U1;
stress1 = D1*strain1;

fprintf('Area A (mm^2):\n');
disp(A1);
fprintf('B matrix (1/mm):\n');
disp(B1);
fprintf('D matrix (MPa):\n');
disp(D1);
fprintf('Element stiffness Ke (N/mm):\n');
disp(Ke1);
fprintf('Nodal displacement vector [u1 v1 u2 v2 u3 v3] (mm):\n');
disp(U1);
fprintf('Strain [epsilon_x epsilon_y gamma_xy]:\n');
disp(strain1);
fprintf('Stress [sigma_x sigma_y tau_xy] (MPa):\n');
disp(stress1);


%% FUNCTIONS

function D = plane_stress_matrix(E, nu)
    D = E/(1-nu^2) * [1 nu 0; nu 1 0; 0 0 (1-nu)/2];
    end
    
    function [A, B, D, Ke] = cst_matrices(coords, E, nu, thickness)
    
    % separate x and y coordinates
    x = coords(:,1);
    y = coords(:,2);
    
    % calculate the area of the triangle
    twoA = x(1)*(y(2)-y(3)) + x(2)*(y(3)-y(1)) + x(3)*(y(1)-y(2));
    A = abs(twoA)/2;
    
    if A <= 0
        error('CST element has zero area.');
    end
    
    % calculate the B matrix
    beta = [y(2)-y(3); y(3)-y(1); y(1)-y(2)];
    gamma = [x(3)-x(2); x(1)-x(3); x(2)-x(1)];
    B = [beta(1) 0 beta(2) 0 beta(3) 0; ...
        0 gamma(1) 0 gamma(2) 0 gamma(3); ...
        gamma(1) beta(1) gamma(2) beta(2) gamma(3) beta(3)]/(2*A);
    
    % calculate the D matrix and element stiffness matrix
    D = plane_stress_matrix(E, nu);
    Ke = thickness*A*(B.'*D*B);
    end

function U = solve_with_constraints(K, F, fixedDofs)
    allDofs = (1:length(F)).';
    freeDofs = setdiff(allDofs, fixedDofs);
    U = zeros(size(F));
    U(freeDofs) = K(freeDofs,freeDofs)\F(freeDofs);
end
\end{lstlisting}

\subsection{Results}
The result obtained from the code, formatted for $\LaTeX$, is given below:
The cross-sectional area of the element is given by:

\[
A = 0.5000 \;\text{mm}^2
\]

The strain-displacement matrix (\(\mathbf{B}\)) is given by:

\[
\mathbf{B} =
\begin{bmatrix}
-1 & 0 & 1 & 0 & 0 & 0 \\
0 & -1 & 0 & 0 & 0 & 1 \\
-1 & -1 & 0 & 1 & 1 & 0
\end{bmatrix}
\;\text{mm}^{-1}
\]

The \(\mathbf{D}\) matrix  is given by:

\[
\mathbf{D} =
10^5
\begin{bmatrix}
2.1978 & 0.6593 & 0 \\
0.6593 & 2.1978 & 0 \\
0 & 0 & 0.7692
\end{bmatrix}
\;\text{MPa}
\]

The element stiffness matrix (\(\mathbf{K}_e\)) is given by:

\[
\mathbf{K}_e =
10^5
\begin{bmatrix}
1.4835 & 0.7143 & -1.0989 & -0.3846 & -0.3846 & -0.3297 \\
0.7143 & 1.4835 & -0.3297 & -0.3846 & -0.3846 & -1.0989 \\
-1.0989 & -0.3297 & 1.0989 & 0 & 0 & 0.3297 \\
-0.3846 & -0.3846 & 0 & 0.3846 & 0.3846 & 0 \\
-0.3846 & -0.3846 & 0 & 0.3846 & 0.3846 & 0 \\
-0.3297 & -1.0989 & 0.3297 & 0 & 0 & 1.0989
\end{bmatrix}
\;\text{N/mm}
\]

The nodal displacement vector, arranged in the order \([u_1, v_1, u_2, v_2, u_3, v_3]^T\), is given by:

\[
\mathbf{d}_e =
10^{-3}
\begin{bmatrix}
0 \\
0 \\
0.3000 \\
0 \\
0 \\
-1.0000
\end{bmatrix}
\;\text{mm}
\]

The corresponding strain vector, consisting of the normal strains and engineering shear strain, is given by:

\[
\boldsymbol{\varepsilon} =
10^{-3}
\begin{bmatrix}
0.3000 \\
-1.0000 \\
0
\end{bmatrix}
\]

The resulting stress vector is given by:

\[
\boldsymbol{\sigma} =
\begin{bmatrix}
0 \\
-200 \\
0
\end{bmatrix}
\;\text{MPa}
\]

\subsection{Discussion}

The deformation is illustrated in Figure~\ref{lab_4_1} below.
\begin{figure}[!h]
    \centering
    \includegraphics[width=0.5\linewidth]{figures/lab_4_1.png}
    \caption{Deformed CST element}
    \label{lab_4_1}
\end{figure}

As you can see, the node 2 (top) moved downward by 0.001 mm, whereas the node . The vertical reaction at node 1 (bottom left) is +100 N, balancing the applied load. The node 3 (bottom right) moved 0.0003 mm to the right, relatively less compared to the node 1 displacement. Because a linear CST has constant $B$, strain and stress are uniform throughout the element. 

The stress is compressive in the vertical direction, which is consistent with the downward force applied at node 3. The shear stress and strain are zero, $\sigma_{xy} = \gamma_{xy} = 0$. This is expected because the triangle is right-angled and the applied load is vertical, resulting in no shear deformation.

\section{Patch Test}

The objective is to perform a patch test on a unit square divided into two triangular CST
elements. We will obtain a linear displacement field at the four nodes, calculate
the strain in each element, and compare the numerical strain with the exact
strain obtained by differentiating the prescribed displacement field. The
test is used to verify that the CST formulation reproduces a compatible
linear displacement field exactly. The system is given as:
\begin{table}[h!]
\centering
\begin{tabular}{ccl}
\toprule
\textbf{Quantity} & \textbf{Property / Value}\\
\midrule
Node 1 & $(0,0)\ \si{\milli\metre}$ \\
Node 2 & $(1,0)\ \si{\milli\metre}$ \\
Node 3 & $(1,1)\ \si{\milli\metre}$ \\
Node 4 & $(0,1)\ \si{\milli\metre}$ \\
Element 1 & Nodes 1-2-3\\
Element 2 & Nodes 1-3-4\\
$u$ & $0.001x$\\
$v$ & $0.002y$\\
\bottomrule
\end{tabular}
\end{table}

\subsection{Code}
\begin{lstlisting}[style=matlabcode]
% setting up a 1x1 square patch with two triangular elements
nodes2 = [0 0; 1 0; 1 1; 0 1];
elements2 = [1 2 3; 1 3 4];
nNodes2 = size(nodes2, 1);
nDofs2 = 2*nNodes2;

% Assemble the global stiffness matrix for the patch test
[K2, elementData2] = assemble_global_stiffness(nodes2, elements2, E, nu, t);

% Prescribed displacement field: u = 0.001*x, v = 0.002*y.
U2 = zeros(nDofs2, 1);
U2(1:2:end) = 0.001*nodes2(:,1); % setting u displacements
U2(2:2:end) = 0.002*nodes2(:,2); % setting v displacements
exactStrain2 = [0.001; 0.002; 0];

fprintf('Prescribed nodal displacement vector (mm):\n');
disp(U2);
fprintf('Element results [element area eps_x eps_y gamma_xy sigma_x sigma_y tau_xy]:\n');

patchResults = zeros(size(elements2,1), 8);
for e = 1:size(elements2,1)

    % select the degrees of freedom for this element
    dofs = element_dofs(elements2(e,:));

    % Extract the displacements for this element
    ue = U2(dofs);

    % Compute the strain and stress for this element
    strain = elementData2(e).B*ue;
    stress = elementData2(e).D*strain;

    % Store the results for this element
    patchResults(e,:) = [e, elementData2(e).area, strain.', stress.'];
    fprintf('%d  %.6f  %.6f  %.6f  %.6f  %.6f  %.6f  %.6f\n', ...
        patchResults(e,:));
end

strainError2 = patchResults(:,3:5) - repmat(exactStrain2.', size(elements2,1), 1);
fprintf('Maximum absolute strain error: %.3e\n', max(abs(strainError2(:))));


%% more LOCAL FUNCTIONS (not repeating the ones from Q1)


function dofs = element_dofs(conn)
    % conn is a vector of node indices for the element
    
    % select the degrees of freedom for the element
    dofs = reshape([2*conn-1; 2*conn], [], 1);
end

function [K, elementData] = assemble_global_stiffness(nodes, elements, E, nu, thickness)
    
    % setting up the number of degrees of freedom for the global stiffness matrix
    nDofs = 2*size(nodes,1);
    
    % initialize the global stiffness matrix and element data structure
    K = sparse(nDofs, nDofs);
    elementData = repmat(struct('area', [], 'B', [], 'D', [], 'Ke', []), ...
        size(elements,1), 1);
    
    for e = 1:size(elements,1)
        % select the nodes for this element
        conn = elements(e,:);
    
        % compute the element matrices using the CST formulation
        [A, B, D, Ke] = cst_matrices(nodes(conn,:), E, nu, thickness);
    
        % assemble the global stiffness matrix and store element data
        dofs = element_dofs(conn);
        K(dofs,dofs) = K(dofs,dofs) + Ke;
    
        % set the element data for this element
        elementData(e).area = A;
        elementData(e).B = B;
        elementData(e).D = D;
        elementData(e).Ke = Ke;
    end
end
\end{lstlisting}

\subsection{Results}
The prescribed nodal displacement vector, arranged in the order of the corresponding nodal degrees of freedom, is given by:

\[
\mathbf{d}_p =
\begin{bmatrix}
0 \\
0 \\
0.0010 \\
0 \\
0.0010 \\
0.0020 \\
0 \\
0.0020
\end{bmatrix}
\;\text{mm}
\]

The element-wise results, including the area, strain components, and stress components, are summarized below:

\[
\begin{array}{c|c|ccc|ccc}
\text{Element} & A\;(\text{mm}^2)
& \varepsilon_x & \varepsilon_y & \gamma_{xy}
& \sigma_x & \sigma_y & \tau_{xy} \\
\hline
1 & 0.5000 & 0.0010 & 0.0020 & 0.0000
& 351.6484 & 505.4945 & 0.0000 \\
2 & 0.5000 & 0.0010 & 0.0020 & 0.0000
& 351.6484 & 505.4945 & 0.0000
\end{array}
\]

The maximum absolute strain error obtained from the analysis is:

\[
\text{Maximum absolute strain error} = 0.000 \times 10^{0}
\]

\subsection{Discussion}
The deformation is illustrated in Figure~\ref{lab_4_2} below.
\begin{figure}[!h]
    \centering
    \includegraphics[width=0.5\linewidth]{figures/lab_4_2.png}
    \caption{Deformed elements from patch test}
    \label{lab_4_2}
\end{figure}

The prescribed nodal values are $(u,v)=[(0,0),(0.001,0),(0.001,0.002),(0,0.002)]$ mm. Both triangular elements recover $\epsilon_x=0.001$, $\epsilon_y=0.002$, and $\gamma_{xy}=0$ to machine precision.

This is a successful patch test: a CST mesh reproduces a compatible linear displacement field exactly.

\section{Cantilever Beam using CST Elements}

The objective is to model a rectangular cantilever beam using triangular CST elements. We will generate a
structured mesh, divide every rectangular mesh cell into two triangles,
assemble the global stiffness matrix, apply the fixed-end boundary
conditions, and calculate the vertical displacement at the loaded tip. We will also repeat the analysis for progressively finer meshes and compare the finite
element tip displacement with the Euler--Bernoulli analytical solution. The system is given as:
\begin{table}[h!]
\centering
\begin{tabular}{ll}
\toprule
\textbf{Quantity} & \textbf{Value} \\
\midrule
Beam length, $L$ & $\SI{1}{\metre}=\SI{1000}{\milli\metre}$ \\
Beam height, $H$ & $\SI{0.1}{\metre}=\SI{100}{\milli\metre}$ \\
Thickness, $t$ & $\SI{1}{\milli\metre}$ \\
Young's modulus, $E$ & $\SI{210000}{\mega\pascal}$ \\
Poisson's ratio, $\nu$ & $0.30$ \\
Applied tip load, $F_y$ & $-\SI{100}{\newton}$ \\
Left-edge boundary condition & $u=v=0$ \\
\bottomrule
\end{tabular}
\end{table}


\subsection{Code}
\begin{lstlisting}[style=matlabcode]
% dimensions 
L = 1000;  % mm
H = 100;   % mm
t = 1;  % mm

% load 
Fy = -100;  % N

% Material properties
E = 200000;  % MPa
nu = 0.30; 
I = t*H^3/12; % Moment of inertia for a rectangular cross-section


% different mesh cases for convergence study
meshCases = [10 2; 20 4; 40 8; 80 16];
nCases = size(meshCases, 1); % no. of mesh cases

% initialize results storage
beamResults = zeros(nCases, 8);
beamSolutions = cell(nCases, 1);

% calculate the analytical tip displacement magnitude using Euler-Bernoulli beam theory
analyticalTipMagnitude = abs(Fy)*L^3/(3*E*I);

for iCase = 1:nCases
    % generate the structured mesh
    Nx = meshCases(iCase,1);
    Ny = meshCases(iCase,2);
    [nodes, elements, nodeId] = structured_mesh(Nx, Ny, L, H);

    % assemble the global stiffness matrix
    [K, ~] = assemble_global_stiffness(nodes, elements, E, nu, t);

    % set no. of nodes and initialize the global force vector
    nNodes = size(nodes,1);
    F = zeros(2*nNodes, 1);


    tipNode = nodeId(Nx, floor(Ny/2)); % find node ID of the tip node at the free end
    tipVDof = 2*tipNode; % vertical degree of freedom for the tip node
    F(tipVDof) = Fy; % apply the vertical load at the tip node

    % fix the nodes at the clamped end
    fixedNodes = zeros(Ny+1, 1);
    for iy = 0:Ny
        fixedNodes(iy+1) = nodeId(0, iy);
    end
    fixedDofs = reshape([2*fixedNodes-1, 2*fixedNodes].', [], 1);

    % solve for the displacements with the applied constraints
    U = solve_with_constraints(K, F, fixedDofs);

    % isolate the tip displacement and compute the relative error
    tipDisplacement = U(tipVDof);
    relativeError = abs(abs(tipDisplacement) - analyticalTipMagnitude) ...
        / analyticalTipMagnitude * 100;

    % store the results for this mesh case
    beamResults(iCase,:) = [Nx, Ny, nNodes, size(elements,1), ...
        tipDisplacement, abs(tipDisplacement), ...
        analyticalTipMagnitude, relativeError];
    beamSolutions{iCase} = struct('nodes', nodes, 'elements', elements, ...
        'U', U, 'tipNode', tipNode);
end

fprintf(['Nx  Ny  Nodes  Elements  FEM tip v (mm)  FEM magnitude (mm)  ', ...
    'Analytical (mm)  Error (%%)\n']);
fprintf('%2d  %2d  %5d  %8d  %16.6f  %19.6f  %16.6f  %8.3f\n', ...
    beamResults.');

fprintf('\nSecond moment of area I = %.6f mm^4\n', I);
fprintf('Analytical tip displacement magnitude = %.6f mm\n', ...
    analyticalTipMagnitude);

%% more local functions (excluding ones from Q1 and Q2 to avoid repetition)


function [nodes, elements, nodeId] = structured_mesh(Nx, Ny, L, H)

    % generate a structured mesh of triangular elements over a rectangular domain
    x = linspace(0, L, Nx+1);
    y = linspace(0, H, Ny+1);
    [X, Y] = meshgrid(x, y);
    
    % create the node coordinates and a function to map (ix, iy) to node ID
    nodes = [X(:), Y(:)];
    nodeId = @(ix, iy) iy + 1 + ix*(Ny+1);
    
    % create the triangular elements by splitting each rectangle into two triangles
    elements = zeros(2*Nx*Ny, 3);
    e = 0;
    for ix = 0:Nx-1
        for iy = 0:Ny-1
            n1 = nodeId(ix, iy);
            n2 = nodeId(ix+1, iy);
            n3 = nodeId(ix+1, iy+1);
            n4 = nodeId(ix, iy+1);
            e = e + 1;
            elements(e,:) = [n1 n2 n3];
            e = e + 1;
            elements(e,:) = [n1 n3 n4];
        end
    end
end
\end{lstlisting}
\subsection{Results}
The second moment of area of the beam cross-section is calculated as:

\[
I =\frac{tH^3}{12}= 83333.33 \;\text{mm}^4
\]

The analytical tip displacement magnitude is given by:

\[
|\delta_{\text{analytical}}| = \frac{|F|L^3}{3EI}= 2.00 \;\text{mm}
\]
The convergence study was conducted using progressively refined meshes with different numbers of elements along the \(x\)- and \(y\)-directions. The resulting finite element tip displacements and percentage errors are summarized in the following table:
\[
\begin{array}{c|c|c|c|c|c|c}
N_x & N_y & \text{Nodes} & \text{Elements}
& \text{FEM tip }v
& \text{Error} \\
& & & & (\text{mm}) & (\%) \\
\hline
10 & 2 & 33 & 40 & -0.742215   & 62.889 \\
20 & 4 & 105 & 160 & -1.403156   & 29.842 \\
40 & 8 & 369 & 640 & -1.813238   & 9.338 \\
80 & 16 & 1377 & 2560 & -1.957955  & 2.102
\end{array}
\]
where the relative error was found by
$$\frac{|\delta| - |\delta_{\text{analytical}}|}{|\delta_{\text{analytical}}|}100\%$$

\subsection{Discussion}
The deformation of the cantilever beam, for 10 segments in $x$ direction and 2 segments in $y$ direction, is illustrated in Figure~\ref{lab_4_3_1} below.
\begin{figure}[!h]
    \centering
    \includegraphics[width=0.65\linewidth]{figures/lab_4_3_1.png}
    \caption{Deformed Cantilever beam}
    \label{lab_4_3_1}
\end{figure}

The Euler-Bernoulli reference value is
$$I=83333.33\;\text{mm}^4,\qquad |\delta_{analytical}|=2.00\;\text{mm}.$$

Refinement from 40 to 2560 CST elements increases the predicted downward displacement toward the analytical value. The results gave us −0.74 mm (62.9\% error) for the coarsest mesh and −1.96 mm (2.1\% error)** for the finest mesh. The convergence to the analytical value is depicted in the Figure~\ref{lab_4_3_2} below.
\begin{figure}[!h]
    \centering
    \includegraphics[width=0.69\linewidth]{figures/lab_4_3_2.png}
    \caption{Convergence of Cantilever beam solution}
    \label{lab_4_3_2}
\end{figure}

Increasing the number of elements improves the accuracy, but convergence is slow. The relative error decreases from 62.9\% to 2.1\% as the mesh is refined when the number of elements is over 64 times the original.

CST elements are too stiff on coarse meshes for bending because each triangle has constant strain and cannot represent the beam’s varying bending strain field. Consequently, bending-dominated problems require refinement and have slower convergence than higher-order elements. The solution is to use higher-order elements or refine the mesh, especially in regions of high bending.


